% §7.2 The Binomial Estimator: Properties
%
% Cross-reference resolution status:
%
% DEFINED here:
%   \label{sec:binomial-estimator}              -- this section
%   \label{prop:binomial-unbiased}              -- unbiasedness proposition
%   \label{prop:binomial-variance}              -- variance proposition
%   \label{prop:binomial-mle}                   -- MLE proposition
%   \label{prop:fisher-information}             -- Fisher information / CRLB
%   \label{prop:binomial-efficiency}            -- CRLB-attainment / efficiency
%
% ALREADY DEFINED (chapter 1):
%   \Cref{sec:sampling-rates}                   -- chapter 1 §1.4
%   \Cref{sec:quantifying-noise}                -- chapter 1 §1.2
%
% ALREADY DEFINED (chapter 7, this plan):
%   \Cref{sec:setup-revisited}                  -- §7.1
%
% FORWARD REFS within chapter 7 (resolve in later Plan 7 tasks):
%   \Cref{sec:confidence-intervals}             -- §7.3 (Task 4)
%
% §7.2 retires item (i) of the four-item §1.4 deferral named in §7.1:
% prove $\hat{\fprate} = \mathrm{FP}/n_0$ is unbiased for $\fprate$,
% derive its variance, identify it as the MLE, and confirm it attains
% the Cramer-Rao lower bound via the Fisher information of the
% Binomial$(n_0, \fprate)$ model. The proofs are short and the
% propositions cumulate; §7.3 then takes up the CI question.

\section{The Binomial Estimator: Properties}
\label{sec:binomial-estimator}

\Cref{sec:setup-revisited} named four items the chapter~1 hook left
deferred. Item~(i) is the present section's target: the binomial
estimator $\hat{\fprate} = \mathrm{FP}/n_0$ stated without proof in
\Cref{sec:sampling-rates} is unbiased for $\fprate$, has the variance
$\fprate(1 - \fprate)/n_0$ claimed there, equals the
maximum-likelihood estimator, and attains the Cramer-Rao lower bound.
Four short propositions, each with a one-paragraph proof, settle
these claims.

\subsection*{Setup}

Fix the truth-$\bot$ slice of the sample. Of the $n$ trials of
\Cref{sec:sampling-rates}, the $n_0$ trials whose truth is $\bot$ are
the only ones a false positive can come from. For each such trial $i
\in \{1, \dots, n_0\}$, let $X_i = 1$ if the noisy bit reports $\top$
and $X_i = 0$ otherwise. Axiom~1 of the Bernoulli model (element-wise
independence) makes the $X_i$ mutually independent and identically
distributed, each Bernoulli with success probability $\fprate$:
\[
  X_i \;\sim\; \mathrm{Bernoulli}(\fprate),
  \qquad
  X_1, \dots, X_{n_0} \text{ independent.}
\]
The false-positive count is their sum,
\[
  \mathrm{FP} \;=\; \sum_{i=1}^{n_0} X_i
  \;\sim\; \mathrm{Binomial}(n_0, \fprate),
\]
recovering the Binomial form already stated in
\Cref{sec:sampling-rates}. The estimator under study is the
row-conditional empirical proportion
\[
  \hat{\fprate} \;=\; \frac{\mathrm{FP}}{n_0}
  \;=\; \frac{1}{n_0} \sum_{i=1}^{n_0} X_i.
\]
The arguments for $\hat{\fnrate} = \mathrm{FN}/n_1$ on the truth-$\top$
slice are formally identical and are not repeated; everything in this
section transfers by renaming $\fprate \to \fnrate$ and $n_0 \to n_1$.

\subsection*{Unbiasedness}

\begin{proposition}[Unbiasedness]
\label{prop:binomial-unbiased}
The empirical estimator $\hat{\fprate} = \mathrm{FP}/n_0$ is unbiased
for $\fprate$:
\[
  \mathbb{E}[\hat{\fprate}] \;=\; \fprate.
\]
\end{proposition}

\begin{proof}
Linearity of expectation, applied to the sum $\mathrm{FP} =
\sum_{i=1}^{n_0} X_i$, gives
\[
  \mathbb{E}[\mathrm{FP}]
  \;=\; \sum_{i=1}^{n_0} \mathbb{E}[X_i]
  \;=\; \sum_{i=1}^{n_0} \fprate
  \;=\; n_0 \fprate,
\]
because each $X_i \sim \mathrm{Bernoulli}(\fprate)$ has expectation
$\fprate$. Dividing by the deterministic constant $n_0$,
\[
  \mathbb{E}[\hat{\fprate}]
  \;=\; \mathbb{E}\!\left[\frac{\mathrm{FP}}{n_0}\right]
  \;=\; \frac{\mathbb{E}[\mathrm{FP}]}{n_0}
  \;=\; \frac{n_0 \fprate}{n_0}
  \;=\; \fprate,
\]
as claimed.
\end{proof}

The bias of the estimator is zero at every $\fprate \in [0, 1]$ and
every sample size $n_0 \geq 1$, not merely in the limit. This is the
first piece of the §1.4 deferral retired.

\subsection*{Variance}

\begin{proposition}[Variance]
\label{prop:binomial-variance}
The empirical estimator $\hat{\fprate} = \mathrm{FP}/n_0$ has variance
\[
  \mathrm{Var}(\hat{\fprate})
  \;=\; \frac{\fprate(1 - \fprate)}{n_0}.
\]
\end{proposition}

\begin{proof}
The textbook variance of a Binomial random variable is
$\mathrm{Var}(\mathrm{Binomial}(n, p)) = n p (1 - p)$; we recover it
from the Bernoulli decomposition for completeness. Each $X_i$ takes
value $1$ with probability $\fprate$ and $0$ with probability $1 -
\fprate$, so
\[
  \mathrm{Var}(X_i)
  \;=\; \mathbb{E}[X_i^2] - \mathbb{E}[X_i]^2
  \;=\; \fprate - \fprate^2
  \;=\; \fprate(1 - \fprate).
\]
By Axiom~1 the $X_i$ are mutually independent, so the variance of
their sum is the sum of their variances,
\[
  \mathrm{Var}(\mathrm{FP})
  \;=\; \sum_{i=1}^{n_0} \mathrm{Var}(X_i)
  \;=\; n_0 \, \fprate(1 - \fprate).
\]
The scaling rule $\mathrm{Var}(c Z) = c^2 \mathrm{Var}(Z)$ for a
constant $c$ then gives
\[
  \mathrm{Var}(\hat{\fprate})
  \;=\; \mathrm{Var}\!\left(\frac{\mathrm{FP}}{n_0}\right)
  \;=\; \frac{1}{n_0^2} \mathrm{Var}(\mathrm{FP})
  \;=\; \frac{n_0 \, \fprate(1 - \fprate)}{n_0^2}
  \;=\; \frac{\fprate(1 - \fprate)}{n_0},
\]
which is the formula stated.
\end{proof}

Two consequences are worth naming. First, the variance shrinks as
$1/n_0$: doubling the truth-$\bot$ sample halves the variance, and
the standard error shrinks as $1/\sqrt{n_0}$. Second, the numerator
$\fprate(1 - \fprate)$ peaks at $\fprate = 1/2$ and vanishes at the
boundaries $\fprate \in \{0, 1\}$. The estimator is hardest to pin
down when the classifier is maximally noisy and easiest when it is
deterministic, a fact that returns in \Cref{sec:confidence-intervals}
as the reason Wald intervals deform near the boundary.

\subsection*{Maximum likelihood}

\begin{proposition}[The MLE is the empirical estimator]
\label{prop:binomial-mle}
The maximum-likelihood estimator of $\fprate$ from the binomial
observation $\mathrm{FP} \sim \mathrm{Binomial}(n_0, \fprate)$ equals
the row-conditional empirical proportion:
\[
  \hat{\fprate}_{\mathrm{MLE}} \;=\; \frac{\mathrm{FP}}{n_0}
  \;=\; \hat{\fprate}.
\]
\end{proposition}

\begin{proof}
The likelihood of the observed count $\mathrm{FP}$ under the model
$\mathrm{Binomial}(n_0, \fprate)$ is
\[
  L(\fprate)
  \;=\; \binom{n_0}{\mathrm{FP}} \, \fprate^{\mathrm{FP}} \,
        (1 - \fprate)^{n_0 - \mathrm{FP}}.
\]
The binomial coefficient does not depend on $\fprate$ and contributes
an additive constant to the log-likelihood:
\[
  \log L(\fprate)
  \;=\; \log \binom{n_0}{\mathrm{FP}}
        + \mathrm{FP} \log \fprate
        + (n_0 - \mathrm{FP}) \log(1 - \fprate).
\]
Differentiating in $\fprate$ over the open interval $(0, 1)$ gives
\[
  \frac{\partial \log L}{\partial \fprate}
  \;=\; \frac{\mathrm{FP}}{\fprate}
        - \frac{n_0 - \mathrm{FP}}{1 - \fprate}.
\]
The first-order condition $\partial \log L / \partial \fprate = 0$
becomes
\[
  \mathrm{FP} \, (1 - \fprate) \;=\; (n_0 - \mathrm{FP}) \, \fprate,
\]
which simplifies to $\mathrm{FP} - \mathrm{FP} \fprate = n_0 \fprate -
\mathrm{FP} \fprate$, then to $\mathrm{FP} = n_0 \fprate$, and finally
to $\fprate = \mathrm{FP}/n_0$. The second derivative
\[
  \frac{\partial^2 \log L}{\partial \fprate^2}
  \;=\; -\frac{\mathrm{FP}}{\fprate^2}
        - \frac{n_0 - \mathrm{FP}}{(1 - \fprate)^2}
\]
is strictly negative on $(0, 1)$ for any $\mathrm{FP} \in \{1, \dots,
n_0 - 1\}$, so the critical point is the unique interior maximum.
The boundary cases $\mathrm{FP} = 0$ and $\mathrm{FP} = n_0$ assign
$\hat{\fprate}_{\mathrm{MLE}} = 0$ and $\hat{\fprate}_{\mathrm{MLE}} =
1$ respectively, again matching $\mathrm{FP}/n_0$. Hence
$\hat{\fprate}_{\mathrm{MLE}} = \mathrm{FP}/n_0 = \hat{\fprate}$.
\end{proof}

The MLE identity says the empirical estimator was not just
convenient: it is what likelihood theory asks for under the binomial
model. Method-of-moments and maximum-likelihood agree on this problem;
in more general settings they need not, as Exercise~9 of this chapter
illustrates.

\subsection*{Fisher information and the Cramer-Rao bound}

\begin{proposition}[Fisher information and Cramer-Rao bound]
\label{prop:fisher-information}
The Fisher information of $\fprate$ in the
$\mathrm{Binomial}(n_0, \fprate)$ model is
\[
  I(\fprate)
  \;=\; \frac{n_0}{\fprate(1 - \fprate)}.
\]
Consequently, every unbiased estimator $T$ of $\fprate$ satisfies the
Cramer-Rao lower bound
\[
  \mathrm{Var}(T) \;\geq\; \frac{1}{I(\fprate)}
  \;=\; \frac{\fprate(1 - \fprate)}{n_0}.
\]
\end{proposition}

\begin{proof}
The Fisher information for a regular parametric model with
log-likelihood $\log L$ is $I(\fprate) = -\mathbb{E}[\partial^2 \log L
/ \partial \fprate^2]$. The second derivative was computed in the
proof of \Cref{prop:binomial-mle}:
\[
  \frac{\partial^2 \log L}{\partial \fprate^2}
  \;=\; -\frac{\mathrm{FP}}{\fprate^2}
        - \frac{n_0 - \mathrm{FP}}{(1 - \fprate)^2}.
\]
Taking expectations under $\mathrm{FP} \sim \mathrm{Binomial}(n_0,
\fprate)$ uses $\mathbb{E}[\mathrm{FP}] = n_0 \fprate$ and
$\mathbb{E}[n_0 - \mathrm{FP}] = n_0(1 - \fprate)$:
\[
  \mathbb{E}\!\left[\frac{\partial^2 \log L}{\partial \fprate^2}\right]
  \;=\; -\frac{n_0 \fprate}{\fprate^2}
        - \frac{n_0(1 - \fprate)}{(1 - \fprate)^2}
  \;=\; -\frac{n_0}{\fprate} - \frac{n_0}{1 - \fprate}.
\]
Combining the two terms over the common denominator $\fprate(1 -
\fprate)$,
\[
  -\mathbb{E}\!\left[\frac{\partial^2 \log L}{\partial \fprate^2}\right]
  \;=\; \frac{n_0(1 - \fprate) + n_0 \fprate}{\fprate(1 - \fprate)}
  \;=\; \frac{n_0}{\fprate(1 - \fprate)}
  \;=\; I(\fprate).
\]
The Binomial family is regular (the support $\{0, 1, \dots, n_0\}$
does not depend on $\fprate$, the log-likelihood is twice
differentiable in $\fprate$ on $(0, 1)$, and the standard
interchange-of-order conditions hold), so the Cramer-Rao inequality
applies and the bound $\mathrm{Var}(T) \geq 1/I(\fprate) = \fprate(1 -
\fprate)/n_0$ follows by direct substitution.
\end{proof}

\subsection*{Efficiency}

\begin{proposition}[Efficiency]
\label{prop:binomial-efficiency}
The empirical estimator $\hat{\fprate} = \mathrm{FP}/n_0$ attains the
Cramer-Rao lower bound and is therefore an efficient unbiased
estimator of $\fprate$.
\end{proposition}

\begin{proof}
\Cref{prop:binomial-unbiased} establishes that $\hat{\fprate}$ is
unbiased. \Cref{prop:binomial-variance} computes
$\mathrm{Var}(\hat{\fprate}) = \fprate(1 - \fprate)/n_0$.
\Cref{prop:fisher-information} computes the Cramer-Rao bound as
$1/I(\fprate) = \fprate(1 - \fprate)/n_0$. The two quantities are
equal, so $\hat{\fprate}$ attains the bound.
\end{proof}

\subsection*{What has been retired}

The four propositions above establish the four claims the chapter~1
hook deferred for item~(i): $\hat{\fprate} = \mathrm{FP}/n_0$ is
unbiased, has variance $\fprate(1 - \fprate)/n_0$, equals the
maximum-likelihood estimator, and attains the Cramer-Rao lower bound
of $\fprate(1 - \fprate)/n_0$. No unbiased estimator of $\fprate$ can
do better in mean-squared error within the binomial model: any
candidate that reports a lower variance is biased and is trading bias
for variance. Item~(i) of the §1.4 deferral is now retired. What
remains is the question of how to turn the point estimate into a
confidence statement that holds at every $n_0$ and every $\fprate$,
not merely asymptotically. \Cref{sec:confidence-intervals} takes up
that question and retires items~(ii) and~(iii).
